\chapter{Building the simulation system}
\label{ch:topology}

Before a simulation can start, the computer needs a complete description of
the system: every atom, its charge, every bond, and the water and salt around
the protein. The Topology Generator (\key{5a}) builds this description for
you. It writes an input file for \term{tLEaP}, the AmberTools program that
builds simulation systems, runs it, and checks the result.

\begin{glance}
\glanceitem{What it does}{Turns your prepared protein into the two files an
  Amber simulation needs, adding a water box and salt on the way.}
\glanceitem{You need it when}{You want to run a simulation (\key{5b}) or
  prepare a QM/MM model from a solvated system.}
\glanceitem{Before you start}{A loaded structure (\key{1a}). Optional:
  cleaned (\key{4a}), protonated (\key{4e}). Disulfides, hemes, metals and
  ligands must be prepared first (see Section~\ref{sec:topology-special}).}
\glanceitem{Where to find it}{\menupath{5a}{Topology Generator}, in section
  \textsf{5.\ SIMULATION SETUP \& EXECUTION}.}
\glanceitem{You get}{\file{<name>_tleap.in}, \file{<name>.prmtop},
  \file{<name>.rst7}, \file{<name>_parmed.pdb}, \file{leap.log}.}
\end{glance}

\section{Topology and coordinates}\label{sec:topology-files}

An Amber simulation reads two files:

\begin{itemize}
  \item The \term{topology} file (\file{.prmtop}, sometimes called
        \file{.parm7}) is the parts list. It names every atom, gives its
        charge and type, and lists every bond, angle and twist together with
        the force-field numbers that describe them. It does not say where
        the atoms are.
  \item The \term{coordinate} file (\file{.rst7}, sometimes \file{.inpcrd})
        gives the position of every atom and the size of the box.
\end{itemize}

The two files belong together: atom number 500 in one is atom number 500 in
the other. ProPrep names them after your structure, for example
\file{1UBQ.prmtop} and \file{1UBQ.rst7}.

\section{The Topology Generator menu}

\begin{screen}
══ Topology Generator Menu ══
  1. ✓ Edit redox site bond definitions
  2. ✓ Generate tLEaP input for single state
  3. ○ Generate tLEaP inputs for all redox microstates
       ⚠ Need batch microstate metadata from Redox Site Prep
  4. ○ Generate prmtop/rst7 files from tLEaP input files
       ⚠ Need to generate tLEaP input first
  5. ○ Refine protonation states via PB (PBSA)
       ⚠ Need to generate topology files first
  6. ○ Generate cpin/cein file for constant pH or redox MD
       ⚠ Need to generate topology files first

→ Suggestion: Edit bonds (option 1) if needed, then generate tLEaP input for
single state (option 2)
\end{screen}

A tick (\texttt{✓}) means the option can be used now, a circle (\texttt{○})
that it is blocked, with the reason on the yellow line below. A dot
(\texttt{●}) marks a step you have already done. Choosing a blocked option
redraws the menu without a message.

The usual route is \key{2} (write the tLEaP input) followed by \key{4} (run
tLEaP). Option~\key{1} is needed only for special bonds, such as disulfides.

\section{Step by step: a solvated protein}\label{sec:topology-steps}

This is the route that worked for ubiquitin (\pdb{1UBQ}). The tutorial in
Chapter~\ref{ch:tutorial} shows every screen; here is the short version.

\begin{steps}
  \item Type \key{5a}, then \key{2}.
  \item Choose the protein force field (Section~\ref{sec:topology-ff}). Press
        \key{Enter} for ff19SB.
  \item At \prompt{Select modified amino acids ()}, press \key{Enter} (none).
  \item Type \key{1} (None) for DNA, RNA, carbohydrate, lipid and small
        molecule force fields, unless your system contains them.
  \item Choose the water model (Section~\ref{sec:topology-water}) and press
        \key{Enter} at the ion parameters.
  \item Choose explicit solvent, the box shape and the buffer
        (Section~\ref{sec:topology-box}).
  \item Choose the salt (Section~\ref{sec:topology-salt}).
  \item ProPrep shows the tLEaP input. Type \key{n} at \prompt{Edit the tLEaP
        template?} unless you want to change it
        (Section~\ref{sec:topology-edit}).
  \item At \prompt{Do you want to continue without bond definitions? [y/n]
        (n)}, type \key{y}.
  \item Press \key{Enter} to accept the file name, and \key{Enter} again to
        view the file.
  \item Type \key{4} and press \key{Enter} at \prompt{Run
        \textless name\textgreater\_tleap.in?}
        (Section~\ref{sec:topology-run}).
\end{steps}

\begin{watchout}
The bond question defaults to \key{n}. In version 1.23.0, \key{n} abandons
the whole step: ProPrep prints \prompt{Please run RedoxSite detection and try
again.} and returns to the menu without writing a file. For a protein without
special bonds, always type \key{y}. Your earlier answers are kept, so a second
attempt is quick.
\end{watchout}

\section{Choosing a force field}\label{sec:topology-ff}

A \term{force field} is the set of rules and numbers that tells the
simulation how strongly atoms attract, repel and hold on to each other. You
choose one for each kind of molecule in the system. Chapter~\ref{ch:forcefields}
explains how to find or make parameters for molecules that no force field
covers.

\begin{tabularx}{\linewidth}{@{}rlL@{}}
\toprule
\textbf{\#} & \textbf{Protein force field} & \textbf{ProPrep's note} \\
\midrule
1 & None & No protein in system \\
2 & \textbf{ff19SB} (default) & Latest Amber protein force field; use with OPC water \\
3 & ff14SB & Widely validated, good for most applications \\
4 & ff15ipq & Requires SPC/Eb water \\
5 & fb15 & Use with TIP3P-FB or TIP4P-FB \\
6 & ff03.r1 & Older alternative \\
7 & ff99SBildn & Legacy \\
8 & Constant pH (ff10) & For constant-pH simulations (Section~\ref{sec:topology-cpin}) \\
9 & Constant Redox (ff10) & For constant-redox simulations \\
10 & Constant pH + Redox (ff10) & Both \\
\bottomrule
\end{tabularx}

\medskip
Use \textbf{ff19SB} unless you have a reason to match earlier work done with
another force field. Pick ff14SB if you want to stay with the most widely
used older choice.

The next questions ask for force fields for other kinds of molecule:

\begin{tabularx}{\linewidth}{@{}lL@{}}
\toprule
\textbf{Question} & \textbf{Choices (default first)} \\
\midrule
Modified amino acids & none (Enter); 2 Modified AA ff19SB; 3 Modified ff14SB;
  4--6 phosphorylated (ff19SB, ff14SB, ff99SB); 7 mimetic and 8 fluorinated
  (ff15ipq). Several may be given, separated by commas. \\
DNA & OL24; 1 None; OL21; OL15; bsc1; tumuc1. Only OL24 and OL21 suit ff19SB. \\
RNA & OL3; 1 None; LJbb; ROC; YIL; Shaw; modrna08 \\
Carbohydrate & GLYCAM\_06j; 1 None; GLYCAM\_06EPb \\
Lipid & lipid21; 1 None; lipid17 \\
Small molecule & GAFF2; 1 None; GAFF \\
\bottomrule
\end{tabularx}

\begin{tip}
For DNA, RNA, carbohydrates, lipids and small molecules the default is a real
force field, not \prompt{None}, even when your system contains none of these.
Extra force fields do no harm to a protein, but type \key{1} for each one you
do not need, so that the input file stays short and clear.
\end{tip}

\section{Choosing a water model}\label{sec:topology-water}

The \term{water model} is the recipe for a single water molecule: where its
charges sit and how big it is. Models differ in accuracy and speed. A
\emph{three-point} model has three charged sites (the three atoms); a
\emph{four-point} model adds one extra site, which makes it more accurate but
a little slower.

\begin{tabularx}{\linewidth}{@{}rlL@{}}
\toprule
\textbf{\#} & \textbf{Water} & \textbf{ProPrep's note} \\
\midrule
1 & None & Implicit solvent, no water molecules (Section~\ref{sec:topology-implicit}) \\
2 & OPC & 4-point; recommended for ff19SB \\
3 & OPC3 & 3-point OPC, faster, good accuracy \\
4 & OPC3-pol & Polarizable OPC3 \\
5 & TIP4P-Ew & 4-point, good alternative \\
6 & TIP3P & 3-point, traditional; recommended for the constant-pH force field \\
7 & SPC/E & 3-point \\
8 & SPC/Eb & Required for ff15ipq \\
9 & TIP5P & 5-point (use with GLYCAM\_06EPb) \\
10, 11 & TIP3P-FB, TIP4P-FB & For fb15 \\
\bottomrule
\end{tabularx}

\medskip
The default is the model recommended for your protein force field: OPC for
ff19SB, TIP3P for the constant-pH force field. When you run \key{2} again, the
default becomes your previous choice. The ubiquitin tutorial used ff19SB with
TIP3P; the crambin and cytochrome~c tests used ff19SB with OPC.

\begin{watchout}
ProPrep does not warn you if you combine a force field with a water model it
does not recommend. The TIP3P line reads ``not recommended'' for ff19SB but
``Recommended'' for the constant-pH force field. Check the
\prompt{Recommended} label next to the water model before you choose.
\end{watchout}

\subsection{Ion parameters}

Next comes \prompt{Select divalent+ ion parameters}. These numbers describe
metal ions such as Mg\textsuperscript{2+}, Ca\textsuperscript{2+} and
Zn\textsuperscript{2+}; sodium and chloride parameters are loaded with the
water model anyway. The default always follows the water model (for example
\prompt{12-6-4 TIP3P} with TIP3P, \prompt{12-6-4 OPC} with OPC). Press
\key{Enter}. With a \term{12-6-4} set, ProPrep later adds a small correction
for ions to the topology with ParmEd (Section~\ref{sec:topology-run}).

\section{Box shape and size}\label{sec:topology-box}

After the force fields, ProPrep asks how to surround the protein.

\begin{screen}
Solvent Model
  1. Explicit solvent (solvation box with water molecules)
  2. Implicit solvent (GB/PB continuum model, no water box)
Select solvent model [1/2] (1): §ans«»

Box Shape
  1. Truncated octahedron (Recommended) - ~29% fewer waters than rectangular
  2. Rectangular box
Select box shape [1/2] (1): §ans«»

Buffer Distance
Distance from solute to box edge. 8-10 A typical, 12-15 A for large
conformational changes.
solvateOct accepts only a single uniform buffer in this AmberTools build.
Switch to a rectangular box if you need per-axis buffers.
Buffer distance (Angstroms) (10.0): §ans«10»
\end{screen}

\begin{itemize}
  \item \textbf{Explicit solvent} puts real water molecules around the
        protein. Choose it for almost all work.
  \item The simulation box repeats itself in every direction, like tiles
        (this is called \term{periodic boundaries}). A \term{truncated
        octahedron} is a box with its corners cut off. For a roughly round
        protein it needs about 29\,\% less water than a rectangular box, so
        the simulation runs faster.
  \item The \term{buffer} is the smallest distance between the protein and
        the edge of the box. 10\,Å is usual; use 12--15\,Å if the protein may
        unfold or change shape a lot. ProPrep accepts 5--25\,Å.
\end{itemize}

A \textbf{rectangular box} (\key{2}) adds a question about the buffer:

\begin{screen}
  1. Uniform (Recommended) - same buffer in x, y, z
  2. Per-axis - different x, y, z buffers
Select buffer mode [1/2] (1): §ans«2»
Buffer in X direction (Angstroms) (10.0): §ans«10»
Buffer in Y direction (Angstroms) (10.0): §ans«10»
Buffer in Z direction (Angstroms) (10.0): §ans«15»
\end{screen}

A per-axis buffer is useful for long, thin molecules. The truncated
octahedron accepts only one value for all directions.

If you type something that is not a number, ProPrep says \prompt{Please enter
a number}; a value outside the range gives \prompt{Value should be between 5.0
and 25.0 Angstroms}. In both cases the question is asked again.

\section{Salt and neutralizing ions}\label{sec:topology-salt}

Two things decide how many ions go into the box:

\begin{itemize}
  \item \textbf{Neutralizing.} The standard way of computing long-range
        electric forces in a periodic box needs a box with zero total charge.
        If the protein is charged, oppositely charged \term{counter-ions} are
        added.
  \item \textbf{Salt concentration.} Real cells contain about 150\,mM salt,
        which screens charges and affects how the protein behaves. Extra pairs
        of ions are added to reach the concentration you choose.
\end{itemize}

For each salt, ProPrep asks for a positive ion (\prompt{Cation}), a negative
ion (\prompt{Anion}) and a concentration:

\begin{tabularx}{\linewidth}{@{}lL@{}}
\toprule
\textbf{Question} & \textbf{Choices (default first)} \\
\midrule
Cation & Na\textsuperscript{+}; K\textsuperscript{+}; Li\textsuperscript{+};
  Cs\textsuperscript{+}; Rb\textsuperscript{+}; Mg\textsuperscript{2+};
  Ca\textsuperscript{2+}; Zn\textsuperscript{2+}. A 2+ ion gives a salt with
  two anions, such as MgCl\textsubscript{2}. \\
Anion & Cl\textsuperscript{--}; Br\textsuperscript{--}; F\textsuperscript{--};
  I\textsuperscript{--} \\
Concentration & 1 150\,mM; 2 100\,mM; 3 50\,mM; 4 Neutralize only (0\,mM);
  5 Custom (0--1000\,mM) \\
Add another salt? & \key{n} (default) or \key{y} \\
\bottomrule
\end{tabularx}

\medskip
ProPrep then runs tLEaP for a moment to count the water molecules and the
charge of your system, and shows how it computed the ion numbers:

\begin{screen}
Running tLEaP info pass for accurate ion counts...
  → Net charge: +0.0
  → Water count: 4792
  ✓ Got required info, terminating tLEaP early
╭──────── Ion Calculation -- 150 mM NaCl, SPLIT Method ────────╮
...
│    Step 2  formula units:  No = Nw · M / 55.5 = 4792 · 0.15 / 55.5 = 13.0 -> 13
...
│  Result: total charge = system (+0) + ions (+0) = 0  ->  13 Na+ + 13 Cl-
...
\end{screen}

\begin{behindscenes}
ProPrep counts ions with the \term{SPLIT method} (Machado and Pantano,
2020). It shares the protein's charge between the two kinds of ion, so the box
is neutral and reaches the target concentration in one step. The panel also
warns that the box shrinks by 5--10\,\% when the simulation first runs at
constant pressure, so a 150\,mM target ends up at about 158--165\,mM. This is
normal practice.
\end{behindscenes}

\subsection{Two or more salts}

Answer \key{y} at \prompt{Add another salt?} to add a second salt. In our
crambin test we added 150\,mM KCl and 10\,mM MgCl\textsubscript{2}. ProPrep
then asks which salt should neutralize the protein's charge:

\begin{screen}
Charge Neutralization
Choose which salt's ions are used to balance the protein charge via tLEaP's
`addions` command.
  1. KCl
  2. MgCl2
Use which salt for neutralization [1/2] (1): §ans«»
\end{screen}

The result for crambin was 11 K\textsuperscript{+}, 1 Mg\textsuperscript{2+}
and 13 Cl\textsuperscript{--}, with the line \prompt{Charge balance verified}.

\subsection{Neutralize only}

Choice \key{4} adds only the counter-ions. For cytochrome~c (\pdb{1HRC}),
which has a charge of +8, ProPrep wrote \texttt{addions mol Cl- 0}, which tells
tLEaP to add as many chloride ions as needed (here eight) near the positive
charges.

\section{Looking at and editing the tLEaP input}\label{sec:topology-edit}

ProPrep turns your answers into a \term{template}: the tLEaP input with
placeholders that it fills in later. It prints the template with line numbers
and asks \prompt{Edit the tLEaP template? [y/n] (y)}.

\begin{screen}
   4 source leaprc.protein.ff19SB
...
   9 source leaprc.water.tip3p
  10 loadamberparams frcmod.ions1lm_1264_tip3p
  11 loadamberparams frcmod.ions234lm_1264_tip3p
  12
  13 # === CUSTOM ATOM TYPES (auto-filled by ProPrep) ===
  14 # ATOM_TYPES_SECTION
...
  20 mol = loadpdb /Users/you/project/ubq/1UBQ.pdb
...
  23 # BOND_DEFINITIONS_SECTION
...
  26 solvateoct mol TIP3PBOX 10.00
...
  30 addionsrand mol Na+ 13 Cl- 13
...
  33 check mol
...
  36 saveamberparm mol 1UBQ.prmtop 1UBQ.rst7
...
Edit the tLEaP template? [y/n] (y):
\end{screen}

In plain words: load the force fields (\texttt{source},
\texttt{loadamberparams}), read your protein (\texttt{loadpdb}), surround it
with water (\texttt{solvateoct}), add ions (\texttt{addionsrand}), check it
(\texttt{check}) and save the two files (\texttt{saveamberparm}). The lines
ending in \texttt{\_SECTION} are filled in by ProPrep: custom atom types and
parameters for special residues, and bond lines from the bond editor.

Answering \key{y} opens the template in a text editor (\texttt{nano}, or the
editor named in your \texttt{EDITOR} setting). Make your change, then press
\key{Ctrl-X}, \key{y} and \key{Enter} to save. ProPrep says
\prompt{Template updated successfully!}. Your changes are kept and used again
the next time you run \key{2}.

\begin{tip}
The most useful edit is the last line, \texttt{saveamberparm}. It sets the
names of the topology and coordinate files. Change them whenever you make a
second, different system from the same PDB file (see the watch-out in
Section~\ref{sec:topology-implicit}).
\end{tip}

\subsection{Running option 2 again}

The second time you choose \key{2}, ProPrep shows your previous settings and
asks whether to change them:

\begin{screen}
Previous settings: protein: {'name': 'ff19SB', 'leaprc': 'leaprc.protein.ff19SB',
'recommended': True, ...}, water: {'name': 'TIP3P', ...}
Solvation: truncated octahedron, 10.0 A buffer, 150 mM NaCl
Reconfigure all settings (forcefields, solvation, template)? [y/n] (n):
\end{screen}

The settings appear in a computer-style list; the last line,
\prompt{Solvation: \ldots}, is the easiest to read. Press \key{Enter} to keep
everything, or type \key{y} to answer all the questions again.

\subsection{File name and viewing}

After the bond question, ProPrep asks \prompt{Enter output filename for tLEaP
input} with the default \file{<pdb name>_tleap.in}, then \prompt{View the
tLEaP input file? [y/n] (y)}. Viewing shows the finished file in a box, with
the \texttt{\_SECTION} lines replaced by comments such as \texttt{\# No custom
atom types needed}.

\section{Running tLEaP}\label{sec:topology-run}

Option~\key{4} runs tLEaP on the input files in your project folder. With one
file, ProPrep shows a status table and asks \prompt{Run
\textless name\textgreater\_tleap.in? [y/n] (y)}. If the topology already
exists it asks \prompt{Topology already exists. Regenerate \ldots? [y/n] (n)}
instead.

With several input files it offers:

\begin{screen}
Run options:
All topology files already exist.
  1. Regenerate all (2 files)
  2. Select specific files
  3. Skip - all topologies already exist
Select option [1/2/3] (3): §ans«2»

Select files to run (comma-separated numbers, 'all', or 'skip'):
   1. ✓ 1UBQ_gb_tleap.in
   2. ✓ 1UBQ_tleap.in

✓ = topology exists, ○ = needs generation
Enter file numbers (e.g., 1,3,5 or 'all' or 'skip') (skip): §ans«1»
\end{screen}

A successful run of \pdb{1UBQ} looked like this (shortened):

\begin{screen}
✓ Inserted 57 TER record(s) in 1UBQ.pdb
Running: tleap -s -f 1UBQ_tleap.in
✓ tLEaP completed successfully

tLEaP Summary: Exiting LEaP: Errors = 0; Warnings = 13; Notes = 1.
...
Running ParmEd validation on 1UBQ.prmtop...
  ✓ Initial validity check passed
  ✓ Molecule definitions are correct
  Applying 12-6-4 LJ potential (water model: TIP3P)...
  ✓ Added LENNARD_JONES_CCOEF (C4 terms) for 12-6-4 potential
...
All expected topology files were generated successfully!
\end{screen}

What happens, in order:

\begin{enumerate}
  \item \textbf{TER records.} tLEaP needs a \term{TER record} (an
        end-of-molecule line) after every separate molecule, such as each
        crystal water. ProPrep adds missing ones. It edits your PDB file in
        place and keeps the original as \file{<name>.pdb.ter_backup}.
  \item \textbf{tLEaP.} The summary line counts errors, warnings and notes.
        \textbf{Errors = 0} is what matters. Warnings such as \prompt{Close
        contact of 1.3 angstroms between nonbonded atoms H and H1} are normal
        for freshly added hydrogens.
  \item \textbf{ParmEd.} \term{ParmEd}, Amber's topology editor, checks the
        topology, makes sure molecules are defined correctly, and adds the
        12-6-4 ion terms if you chose such a set. It then overwrites the
        \file{.prmtop} with the corrected version.
  \item \textbf{Validation.} A table shows whether the expected
        \file{.prmtop} and \file{.rst7} were found.
\end{enumerate}

\begin{tabularx}{\linewidth}{@{}lL@{}}
\toprule
\textbf{File} & \textbf{Contents} \\
\midrule
\file{<name>.prmtop} & The topology (for ubiquitin 2.9\,MB, 15,729 atoms) \\
\file{<name>.rst7} & The coordinates (574\,kB) \\
\file{<name>_parmed.pdb} & The whole system as a PDB, same atom order as the topology \\
\file{<name>_parmed_replay.py} & A script to repeat the ParmEd step by hand \\
\file{leap.log} & tLEaP's complete report; the place to look when something fails \\
\file{<name>.pdb.ter_backup} & Your PDB as it was before TER records were added \\
\bottomrule
\end{tabularx}

\medskip
The output files are named \file{.prmtop}, not \file{.parm7}; the name comes
from the \texttt{saveamberparm} line of the input. The whole step takes a few
seconds for a small protein.

\section{Other choices}

\subsection{Implicit solvent}\label{sec:topology-implicit}

In \term{implicit solvent} the water is not made of molecules but replaced by
a smooth surrounding medium, the \term{Generalized Born} (GB) model. It is
much cheaper but less realistic. ProPrep switches to it when you choose water
model \key{1} (None) or solvent model \key{2}. The ion question then defaults
to \prompt{13 Default only}, and ProPrep writes no water box.

\begin{watchout}
Our implicit-solvent test on raw \pdb{1UBQ} (ff14SB, no water) failed in
version 1.23.0:
\begin{itemize}
  \item The 58 crystal waters were still in the PDB, and without a water
        model tLEaP did not recognize them: \prompt{tLEaP failed (exit code:
        31)}, \prompt{Unknown residue: HOH}. Remove all waters with
        \key{4a} (Chapter~\ref{ch:cleaning}) before an implicit-solvent build.
  \item ProPrep did not add the radii that a GB simulation needs (the
        \texttt{set default PBradii} line). Add such a line yourself when you
        edit the template.
  \item The new input used the same output names (\file{1UBQ.prmtop},
        \file{1UBQ.rst7}) as the earlier explicit-solvent input. The failed
        run left \file{1UBQ.prmtop} at 0 bytes and so destroyed the good
        topology. Give every variant its own names on the
        \texttt{saveamberparm} line.
\end{itemize}
\end{watchout}

\subsection{Disulfide bonds}\label{sec:topology-disulfide}

A \term{disulfide bond} links the sulfur atoms of two cysteines. tLEaP makes
it only when \emph{both} of these are true:

\begin{enumerate}
  \item Both cysteines are named \textbf{CYX} in the PDB file (the Amber name
        for a cysteine in a disulfide; Appendix~\ref{app:residues}).
  \item The tLEaP input contains a \texttt{bond} line for each pair, for
        example \texttt{bond mol.3.SG mol.40.SG}.
\end{enumerate}

Our tests on crambin (\pdb{1CRN}, three disulfides) showed what happens
otherwise:

\begin{tabularx}{\linewidth}{@{}lL@{}}
\toprule
\textbf{What was done} & \textbf{Result} \\
\midrule
Names left as CYS, no bond lines &
  tLEaP: \prompt{Errors = 42}, \prompt{Could not find bond parameter for atom
  types: SH - SH}. No topology. \\
Names left as CYS, bond lines added &
  The same errors (60). No topology. \\
Names CYX, no bond lines &
  \prompt{Errors = 0}, but the topology has \emph{no} disulfides. Only a
  yellow ParmEd line hints at it. \\
Names CYX \emph{and} bond lines & Correct topology, no warnings. \\
\bottomrule
\end{tabularx}

\medskip
The Protonation State Analyzer (\key{4e}) did \emph{not} rename the cysteines
to CYX. The Redox Site Detector (\key{2a}) followed by the Redox Site Preparer
(\key{4d}) with its disulfide option did (Chapter~\ref{ch:redox}).

\begin{watchout}
A CYX build without bond lines reports success. Look for this line in the
ParmEd output; it means the disulfides are missing:
\begin{screen}
• Detected CYX residue with a Sulfur atom not bonded to another CYX
  Sulfur! Did you forget to make the disulfide bond?
\end{screen}
\end{watchout}

\subsubsection{Adding bond lines with the bond editor}

Option~\key{1} opens the \textsf{tLEaP Bond Editor Dashboard}. If ProPrep has
found no bonds itself it first asks \prompt{Do you want to continue without
bond definitions?}; either answer opens the editor.

\begin{screen}
═══ tLEaP Bond Editor Dashboard ═══
Showing: All Bond Categories
No bonds to display
Commands:
Filters: [d]isulfide [m]etal [p]eptide [o]ther [all] categories
Actions: [a]dd bond, [#] edit bond, [x#] delete bond, [s]ave, [q]uit, [h]elp
Command (q): §ans«a»

Add New Bond
Enter tLEaP bond command (e.g. ...): §ans«bond mol.3.SG mol.40.SG»
Select category:
  1. Covalent Bonds
  2. Disulfide Bonds (SG-SG)
...
  6. Other Bonds
Category (6): §ans«2»
\end{screen}

\begin{steps}
  \item Type \key{a}, then the bond line. Numbers are residue numbers as tLEaP
        sees them (counted from 1 along the file); \texttt{SG} is the sulfur
        atom.
  \item Type \key{2} (Disulfide Bonds). Press \key{Enter} at \prompt{Press
        Enter to continue}.
  \item Repeat for each pair, then type \key{s} to save. ProPrep reports, for
        example, \prompt{Saved 3 bond definitions}.
  \item Run \key{2} again (keep the settings) and then \key{4}. The input
        now contains the lines under \texttt{\# Disulfide bonds (CYS/CYX
        SG-SG)}.
\end{steps}

Other keys: a number edits that bond, \texttt{x} followed by a number deletes
it, \key{d} \key{m} \key{p} \key{o} show one category, \key{h} shows help.
Typing \key{q} \emph{saves} and leaves; there is no ``quit without saving''.

\begin{watchout}
The editor accepts any text as a bond. In our test a stray \key{n} typed at
a prompt was saved as a ``bond'' called \texttt{n}. After each \key{a}, the
pause \prompt{Press Enter to continue} swallows whatever you type. Check the
list before you save.
\end{watchout}

\subsection{Hemes, metals and ligands}\label{sec:topology-special}

Standard force fields know only standard residues. A heme, a metal site or a
drug molecule needs its own names and parameters, which ProPrep prepares with
the Redox Site Detector (\key{2a}), the Redox Site Preparer (\key{4d}) and,
if needed, the Force Field Parameterizer (\key{2c}). The Topology Generator
then fills the \texttt{ATOM\_TYPES} and \texttt{FORCEFIELD\_PARAMETERS}
sections of the template for you.

Going straight to \key{5a} does not work. Cytochrome~c (\pdb{1HRC}) taken
directly from the PDB gave \prompt{tLEaP failed (exit code: 31)} with
\prompt{Errors = 159} and \prompt{Unknown residue: HEC}. ProPrep did not warn
about the heme before running tLEaP. Prepare hemes as described in
Chapter~\ref{ch:redox}, and ligands and metals as in
Chapter~\ref{ch:forcefields}.

\subsection{Many redox microstates at once (option 3)}

When the Redox Site Preparer (\key{4d}) has produced a batch of
\term{microstates} (the same protein with different sites oxidized or
reduced), option~\key{3} writes one tLEaP input for each. Option~\key{2} is
then blocked with \prompt{Batch microstates already generated}.

\begin{watchout}
In version 1.23.0, option~\key{3} silently writes nothing when no site needs
force-field files, for example a batch of crambin microstates that differ only
in their disulfides. ProPrep showed \prompt{Found 4 microstates to process}
and then the menu again, with no input files and no message. Use option~\key{2}
on each microstate PDB instead.
\end{watchout}

\subsection{Starting from an existing topology and coordinates}

If you already have a \file{.prmtop} and \file{.rst7} pair, for example from a
colleague, you can skip this chapter. In the Structure Loader type \key{1a},
\key{2} (Load AMBER topology \& coordinate files), pick the topology in the
file browser (type a folder path to jump there), answer \key{r} at
\prompt{Coordinate source: [r]st7/inpcrd or [t]rajectory}, and pick the
coordinates. ProPrep confirms \prompt{Loaded topology + rst7 pair into
workspace}.

The pair can be used in the Molecular Dynamics Manager (\key{5b}) and the
QM/MM Preparator. Inside \key{5a}, options 2--6 stay blocked, so PB refinement
and cpin files need a topology built in \key{5a}.

\begin{tip}
After you pick the topology, the file browser returns to your project folder,
not the topology's folder. Type the full path of the coordinate file, or
navigate back to it.
\end{tip}

\subsection{Refining protonation states with PB (option 5)}

Option~\key{5} opens \textsf{PB Titrate}, a checklist of eleven steps that
predicts the protonation state of each titratable residue from
\term{Poisson--Boltzmann} electrostatics and can write the result into the
topology. Its defaults are an internal dielectric of 4, 150\,mM ionic
strength, a 0.5\,Å grid and pH~7.0. We opened the checklist and its
information pages but did not run the calculations for this manual.
Opening it creates a folder \file{pb_titrate_work/} in the project. To leave,
type \key{q} and then \key{y} (\key{Enter} alone means ``no'' and keeps you
in the checklist).

\subsection{Files for constant-pH simulations (option 6)}\label{sec:topology-cpin}

In a \term{constant-pH} simulation, residues such as Asp, Glu, His, Lys and
Tyr can gain and lose protons while the simulation runs. Amber needs a
\term{cpin} file listing these residues. Option~\key{6} makes it with the
AmberTools program \texttt{cpinutil.py}.

\begin{steps}
  \item Type \key{6}. Choose \key{1} (Implicit solvent) or \key{2} (Explicit
        solvent).
  \item ProPrep lists the titratable residues it found, with default pKa
        values. Choose which to include: \key{1} all, \key{2} by type,
        \key{3} by number, \key{4} exclude some, \key{5} by pKa range.
  \item Choose the starting states: \key{1} cpinutil defaults, or \key{2} by
        target pH (ProPrep then shows, for example, \prompt{TYR (pKa 9.6): pH
        7.0 \textless{} pKa \textrightarrow{} State 0 = OH (protonated, neutral)}).
  \item Accept the file name.
\end{steps}

For explicit solvent you get \file{<name>.cpin}, a modified topology
\file{<name>_cpin.prmtop} and a script \file{<name>_cpin_generate.sh}.

\begin{watchout}
\begin{itemize}
  \item Build the topology with a \textbf{Constant pH (ff10)} force field
        (\key{8}). With ordinary ff19SB, ProPrep still made a cpin file for
        ubiquitin, but it contained only lysines and a tyrosine: aspartates,
        glutamates and histidines cannot titrate in that topology. ProPrep
        gave no warning.
  \item Even with the constant-pH force field, aspartate and glutamate are
        only found if they were renamed AS4 and GL4 first (for example with
        \key{4e} and its constant-pH option, Chapter~\ref{ch:protonation}).
  \item Making a cpin file switches the Molecular Dynamics Manager to the
        \file{_cpin.prmtop} topology and asks to enable constant-pH MD by
        default. Remove that pair in the first setup step if you want an
        ordinary simulation (Chapter~\ref{ch:runningmd}).
\end{itemize}
\end{watchout}

\section{If something goes wrong}\label{sec:topology-problems}

\subsection{ProPrep says it worked, but it did not}

In version 1.23.0 several messages can be too optimistic after a failed
tLEaP run:

\begin{itemize}
  \item \prompt{tLEaP completed successfully} followed by \prompt{Errors = 42}.
  \item \prompt{Results: 1/1 files processed successfully} together with
        \prompt{Some expected files are missing}.
  \item The menu suggestion \prompt{Topology files generated (1 prmtop
        files)} although the \file{.prmtop} is empty.
  \item A status table showing \texttt{prmtop ✓} for a 0-byte file.
\end{itemize}

Trust the \prompt{Errors =} count and the size of the file. A real topology
for even a small protein is at least several hundred kilobytes. To check the
sizes, open the project folder in your file browser, or type in the Terminal
(inside the project folder):

\begin{terminal}
ls -l *.prmtop *.rst7
\end{terminal}

A line with \texttt{0} before the date means the file is empty.

\subsection{The real error is only in leap.log}

ProPrep shows tLEaP's warnings and the error count, but not always the line
that explains a failure. In the implicit-solvent test the screen showed only
\prompt{Unknown residue: HOH}; the real reason was in \file{leap.log}:

\begin{screen}
FATAL: Atom .R<HOH 120>.A<O 1> does not have a type
\end{screen}

Open \file{leap.log} in the project folder with any text editor and search for
\texttt{FATAL} or \texttt{Error}, or type in the Terminal:

\begin{terminal}
grep -n "FATAL\|Error" leap.log
\end{terminal}

\subsection{Common problems}

\begin{tabularx}{\linewidth}{@{}lL@{}}
\toprule
\textbf{What you see} & \textbf{Why, and what to do} \\
\midrule
Back at the menu, no \file{_tleap.in} written &
  You answered \key{n} (the default) to \prompt{continue without bond
  definitions?}. Run \key{2} again and type \key{y}. \\
\prompt{Unknown residue: HOH} &
  Crystal waters in an implicit-solvent build. Remove waters with \key{4a}. \\
\prompt{Unknown residue: HEC} (or another name) &
  A heme, ligand or metal without parameters. Prepare it first
  (Chapters~\ref{ch:redox} and~\ref{ch:forcefields}). \\
\prompt{Could not find bond parameter for atom types: SH - SH} &
  Disulfide cysteines still named CYS. Rename them to CYX and add bond lines
  (Section~\ref{sec:topology-disulfide}). \\
\prompt{Illegal CONECT record in pdb file} &
  The PDB's CONECT lines no longer match after renaming. Harmless by itself,
  but check that disulfides were made. \\
A good topology became 0 bytes &
  Two input files write the same output names. Give each its own names on the
  \texttt{saveamberparm} line. \\
\prompt{ParmEd validation failed: Could not identify file format} &
  tLEaP did not write a topology. Read \file{leap.log}. \\
\bottomrule
\end{tabularx}

\subsection{ProPrep uses another Amber installation}

Every tLEaP message shows the program's full path, for example\\
\texttt{.../envs/<name>/bin/teLeap}. If that folder is not ProPrep's own, a
different Amber installation is being used. This happens when the
\texttt{AMBERHOME} setting points elsewhere; ProPrep keeps it without a
message. Close ProPrep, type \texttt{unset AMBERHOME} in the Terminal (after
starting the ProPrep environment), and start ProPrep again. A repeat of the
ubiquitin build with only ProPrep's own AmberTools gave the same result
(\prompt{Errors = 0; Warnings = 13}). See also Chapter~\ref{ch:installation}.

More symptoms are listed in Chapter~\ref{ch:troubleshooting}, and the
limitations of version 1.23.0 in Chapter~\ref{ch:knownissues}.
